\begin{claim}{D.2.43}{Proved}{The exact quadratic obstruction and its finite-model survival}
Work at the prime two and put
\[
 D_2E=(E\wedge E)_{hC_2},\qquad L_2E=\Sigma^{-1}D_2E,
 \qquad M_q=M(\mathbb Z/q,7).
\]
For a space in the argument of $D_2$ or $L_2$, its suspension
spectrum is understood. Let
\[
 K=(M_4\vee S^8_r)\cup_{i\eta_7+2r}e^9_s,
 \qquad m_0:M_2\xrightarrow{c_{2,4}}M_4\longrightarrow K,
 \tag{D.2.43a}
\]
where the first map has bottom degree two and top degree one.
Then
\[
 \begin{split}
  \pi_{14}L_2M_2&=\mathbb Z/4\{c_2\},\\
  \pi_{14}L_2K&=\mathbb Z/8\{c\}\oplus\mathbb Z/2\{u\},
  \qquad (L_2m_0)_*c_2=2c.
 \end{split}
 \tag{D.2.43b}
\]
The same target quadratic group applies to
$Z'=(\Sigma^2K(\mathbb F_2,2))\langle7\rangle$.
For the actual Moore map $h_M:M_2\to Z'$ fixed in D.2.41,
the natural metastable boundary satisfies
\[
 h_Ml_4=\partial_{Z'}(\pm2c).
 \tag{D.2.43c}
\]
Consequently the precise remaining nullity test is
\[
 h_Ml_4=0
 \quad\Longleftrightarrow\quad
 2c\in\operatorname{im}\left(
  JH_*:\pi_{15}^{s}Z'\longrightarrow
       \mathbb Z/8\{c\}\oplus\mathbb Z/2\{u\}
 \right),
 \tag{D.2.43d}
\]
where $JH$ is the stable James--Hopf map.

For the finite model $K$ itself, the James--Hopf homomorphism
is zero. Its quadratic boundary is therefore injective, and
\[
 |\partial_K(2c)|=4.
 \tag{D.2.43e}
\]
This class retains exact order four after adding the free
nine-sphere of the target model and any cells of dimensions
ten and eleven. The first possible new contribution to the
James--Hopf image comes from a twelve-cell.
No equality between the finite-model image and the image for
the full $Z'$ is asserted. Thus
\textup{(D.2.43d)} is an exact unresolved target calculation;
it does not settle the fourth-suspension extension problem.
\end{claim}
\begin{proof}
\textbf{The quadratic layer and its integral coordinates.}
The interchange in $D_2$ is the usual symmetric-monoidal
interchange. There is no additional sign representation:
\cite[Section~8, p.~704]{FGHGoodwillie2003} identifies the
second derivative of the identity with $S^{-1}$ carrying
trivial $C_2$-action.

\medskip\noindent
\textbf{The finite quadratic calculation, including the eta attachment.}
Write $D_2E=(E\wedge E)_{hC_2}$ and $L_2E=\Sigma^{-1}D_2E$,
with the ordinary interchange action.  Put $M_q=M(\ZZ/q,7)$,
where $q$ is even.  Its integral label chains have generators
$a_7,b_8$ and $\partial b=qa$.  Use the free $C_2$ resolution
with the convention
\[
 \partial(e_i xy)=e_{i-1}\bigl(xy+(-1)^i\tau(xy)\bigr)
                   +(-1)^ie_i\partial(xy).
 \tag{D.2.43q1}
\]
Thus $\tau(aa)=-aa$, $\tau(ab)=ba$, and $\tau(bb)=bb$.
The cells relevant to integral homology through degree fifteen are
\[
\begin{array}{c|l}
14&P=e_0aa\\
15&x=e_0ab,\quad y=e_0ba,\quad z=e_1aa\\
16&T=e_0bb,\quad h=e_1ab,\quad h'=e_1ba,\quad U=e_2aa.
\end{array}
\]
Their boundaries are
\[
\begin{aligned}
 \partial x&=-qP,&\partial y&=qP,&\partial z&=2P,\\
 \partial T&=q(x+y),&\partial h&=x-y+qz,
       &\partial h'&=-\partial h,\qquad\partial U=0.
\end{aligned}
\tag{D.2.43q2}
\]
Cancel the unit incidence in $\partial h$, so that $y=x+qz$,
and set $c_q=x+(q/2)z$.  Then $c_q$ is an integral cycle and
$\partial T=2q c_q$.  It follows that
\[
 H_{14}(D_2M_q;\ZZ)=\ZZ/2\{P\},\qquad
 H_{15}(D_2M_q;\ZZ)=\ZZ/(2q)\{c_q\}.
 \tag{D.2.43q3}
\]
This chain computation applies to spectra: extension of scalars
to $H\ZZ$ is symmetric monoidal, so the derived coinvariant
tensor complex computes the integral homology of the
homotopy-orbit square.  It does not discard the homotopy
information in eta attachments; that information is checked next.

The bottom inclusion gives $D_2S^7\to D_2M_q$.  Its first three
cells are $P_{14},z_{15},U_{16}$, with $\partial z=2P$ and
$\partial U=0$.  The extended-power basis and Nishida relation
in \cite[\S1.4, Lemma~1.4.1 and the following displayed relation,
p.~4]{EH5Behrens2012} identify these mod-two classes as
$Q^7a,Q^8a,Q^9a$ and give
\[
 \operatorname{Sq}_*^2U
  =\operatorname{Sq}_*^2Q^9a
  =\binom72 Q^7a=P.
 \tag{D.2.43q4}
\]
The fifteen-skeleton of $D_2S^7$ is $M(\ZZ/2,14)$;
its fifteenth stable homotopy group is generated by the bottom
eta.  Hence the sixteen-cell $U$ attaches by, and kills, $P\eta$.
The same class remains zero after mapping to any of the quadratic
spectra considered below.

To see that this accounts for the entire possible Hurewicz kernel,
let $E$ be a spectrum with cells in dimensions at least fourteen
and $H_{14}E=\ZZ/2\{P\}$.  The fiber of
$E\to H\ZZ\wedge E$ is $E\wedge\tau_{\ge1}\mathbb S$.
Its lowest possible homotopy group is
\[
 \pi_{15}(E\wedge\tau_{\ge1}\mathbb S)
       =H_{14}E\otimes\pi_1\mathbb S=\ZZ/2,
\]
and its map into $\pi_{15}E$ is the bottom eta action.
Thus the degree-fifteen stable Hurewicz map is surjective with
kernel generated at most by $P\eta$.  Equation~\textup{(D.2.43q4)}
kills that kernel here, proving the natural isomorphisms
\[
 \pi_{15}D_2M_q\xrightarrow{\ \cong\ }
       H_{15}(D_2M_q;\ZZ)=\ZZ/(2q)\{c_q\}.
 \tag{D.2.43q5}
\]
Equivalently, stable unit cancellation gives the relevant
Chang block with attaching matrix $(2,\eta_{14})$, together
with a Moore $2q$ block whose bottom dimension is fifteen.
Any eta component on that Moore block's top attachment is
absorbed by $U$; no hidden homotopy extension is being inferred
from integral chains alone.

The specified comparison $c_{2,4}:M_2\to M_4$ acts on label chains
by $a\mapsto2a$, $b\mapsto b$.  Therefore
\[
 x\longmapsto2x,\qquad z\longmapsto4z,
 \qquad c_2=x+z\longmapsto2x+4z=2c_4.
 \tag{D.2.43q6}
\]
Naturality of the Hurewicz isomorphism computes the map on
$\pi_{15}$ as $\ZZ/4\to\ZZ/8$, $1\mapsto2$.

Now put $K_8=M_4\vee S^8_r$.  Distributing the two smash factors
and taking homotopy orbits gives the actual stable splitting
\[
 D_2K_8\simeq D_2M_4\vee(M_4\wedge S^8_r)\vee D_2S^8_r.
\]
The middle summand is $M(\ZZ/4,15)$, with bottom class $u=ar$;
the last summand starts in dimension sixteen.  Thus
\[
 \pi_{15}D_2K_8=\ZZ/8\{c_4\}\oplus\ZZ/4\{u\}.
 \tag{D.2.43q7}
\]
Attach the nine-cell of
$K=K_8\cup_{i\eta_7+2r}e^9_s$ and filter its square by the
number of factors in the new cell.  The one-new-factor quotient
is $K_8\wedge S^9$; the two-new-factor quotient is $D_2S^9$,
which starts in dimension eighteen.  Through dimension sixteen,
the former quotient has just the cell $as$.  Its attachment is
the orbit image of $a\wedge(i\eta_7+2r)$, hence
$P\eta+2u$.  There is no extra factor of two on the eta term:
the off-diagonal cross term is an induced free $C_2$-orbit,
whose homotopy orbits give one copy, as in the preceding wedge
splitting.  By \textup{(D.2.43q4)} this imposes exactly $2u=0$.
All further cells have dimension at least seventeen and do not
affect $\pi_{15}$.  Consequently
\[
 \begin{aligned}
 \pi_{14}L_2M_2&=\ZZ/4\{c_2\},\\
 \pi_{14}L_2K&=\ZZ/8\{c_4\}\oplus\ZZ/2\{u\},
       & (L_2m_0)_*c_2&=2c_4.
 \end{aligned}
 \tag{D.2.43q8}
\]

Finally, the extra free nine-sphere in the complete target model
adds a free sixteen-cell to the quadratic spectrum, with zero
attachment, and higher-dimensional cells.  It has no effect on
\textup{(D.2.43q8)}.  A target cell of dimension at least ten
first contributes a cross term in dimension at least seventeen,
because the bottom target cell has dimension seven.  Homotopy
orbits add only nonnegative cell dimensions.  Thus the complete
target's quadratic spectrum has the same fifteenth homotopy
group and the same specified map.  This truncation statement
concerns the quadratic layer; it does not truncate the domain
of the target James--Hopf homomorphism.


\medskip\noindent
\textbf{The actual metastable boundary.}
Theorem~4.2.1 of \cite{EH5Behrens2012}, printed p.~37 in
the opened author version, gives the natural fiber sequence
\[
 P_2\operatorname{Id}(X)\longrightarrow QX
       \xrightarrow{JH}\Omega^\infty D_2\Sigma^\infty X.
 \tag{D.2.43f}
\]
Here $JH$ is adjoint to the weight-two projection of the
Snaith splitting. The following comparison statement in that
source gives
$\pi_tX\cong\pi_tP_2\operatorname{Id}(X)$ for
$t\leq3\operatorname{conn}(X)$. All our spaces in this
paragraph are six-connected, so degrees fourteen and fifteen
are within the stated range. The resulting exact sequence is
\[
 \pi_{15}^{s}X\xrightarrow{JH_*}\pi_{15}D_2X
       \xrightarrow{\partial_X}\pi_{14}X
       \longrightarrow\pi_{14}^{s}X.
 \tag{D.2.43g}
\]

D.2.41, using
\cite[Lemmas~2.7--2.8(ii), p.~523]{EH5MorisugiMukai2000},
identifies the exact stable kernel of $M_2$ as
$\mathbb Z/4\{l_4\}$. The quadratic source in
\textup{(D.2.43b)} also has order four. Hence its boundary
onto that kernel is an isomorphism:
\[
 \partial_{M_2}:\pi_{15}D_2M_2
       \xrightarrow{\ \cong\ }
       \ker(\pi_{14}M_2\to\pi_{14}^{s}M_2).
 \tag{D.2.43h}
\]
In particular $JH_*:\pi_{15}^{s}M_2\to\pi_{15}D_2M_2$
is zero. The generator lifting $l_4$ is unique, up to its
sign convention; there is no unrecorded choice of a source
connecting class.

The stable nine-skeleton of $Z'$ is $K\vee S^9_R$,
by twice suspending the model in \textup{(D.2.22t2)}.
In this low range the model and its specified Moore map
are realized as ordinary maps: their source dimensions are
at most nine and the targets are six-connected, so
Freudenthal gives the required bijection with stable maps.
The pure-$k$ Moore restriction of D.2.41 is the
bottom-two/top-one comparison in \textup{(D.2.43a)}.
The quadratic truncation proved above therefore identifies
its quadratic map with \textup{(D.2.43b)}.
Naturality of \textup{(D.2.43g)--(D.2.43h)} gives
\textup{(D.2.43c)}; exactness gives \textup{(D.2.43d)}.
Unlike the quadratic group, the James--Hopf image cannot
be computed by merely replacing its domain with the
nine-skeleton.

\medskip\noindent
\textbf{Exact order four in the finite model.}
The intrinsic Thom-cell calculation used above also gives
\[
 \pi_{15}D_2S^7=0:
 \tag{D.2.43i}
\]
its integral homology in this degree is zero, and the
sixteen-cell kills its only possible Hurewicz-kernel class,
the bottom eta. By naturality every class induced from
the bottom sphere has zero James--Hopf image.

We first specify an actual boundary calculation.
Let $L=M_4$ or $K$, put $B=L/S^7$, and let
$F=\operatorname{hofib}(L\to B)$. The defining homotopy
pushout square has initial maps $S^7\to L$ and $S^7\to *$,
both seven-connected. Blakers--Massey, in the form
\cite[Section~2, p.~309, and Theorem~2.3, p.~310]{FGHGoodwillie1992},
makes the canonical map $i:S^7\to F$ thirteen-connected.
In both cases $B=\Sigma B'$ with $\dim B'\leq8$. The
fibration boundary restricted to the suspension unit,
$B'\to\Omega B\to F$, therefore factors uniquely through $i$.

This factor is determined by stabilization, since maps
from $B'$ into $S^7$ are in the stable range. To identify it,
compare with the stable fiber:
\[
 F\longrightarrow\operatorname{hofib}(QL\to QB)
       \simeq QS^7.
\]
The map of fiber sequences identifies the stabilized factor
with the desuspension of the stable cofiber connecting map.
The resulting \emph{actual} factors are
\[
 \begin{array}{c|c|c}
 L&B'&B'\longrightarrow S^7\\ \hline
 M_4&S^7&4\iota_7\\
 K&S^7_b\vee M(\mathbb Z/2,7)_r&(4\iota_7,\eta_7q).
 \end{array}
 \tag{D.2.43j}
\]
For the second row, collapsing the bottom gives
\[
 K/S^7\simeq S^8_b\vee M(\mathbb Z/2,8)_r,\qquad
 \delta=(4,\eta_8q):K/S^7\longrightarrow S^8.
 \tag{D.2.43k}
\]
Thus the eta component of the nine-cell is retained in both
the stable connecting map and its actual low-unit factor.
This argument uses a homotopy fiber, not an unstable
covariant exact sequence for a cofiber.

Next, \cite[Introduction, p.~515]{EH5MorisugiMukai2000}
states
\[
 \pi_{14}S^7_{(2)}=\mathbb Z/8\{\sigma'\},\qquad
 [\iota_8,\iota_8]=\pm(2\sigma_8-E\sigma').
 \tag{D.2.43l}
\]
Consequently $E^\infty\sigma'=2\sigma$, up to sign.
The actual class $2E\sigma'\in\pi_{15}S^8$ has stable
image $4\sigma$. Its boundary in the first case of
\textup{(D.2.43j)} is
\[
 i\bigl((4\iota_7)\circ2\sigma'\bigr)=8i\sigma'=0.
 \tag{D.2.43m}
\]
There is no omitted Hilton term: $S^7$ is an H-space,
so its degree-four map induces multiplication by four
on homotopy groups. Equivalently the Hilton correction
vanishes because $[\iota_7,\iota_7]=0$; see also
\cite[Lemma~1.1]{EH5MorisugiMukai2000}.

The actual fibration exact sequence therefore supplies
$\widehat a\in\pi_{15}M_4$ with pinch $2E\sigma'$.
Its stable class has zero James--Hopf image and pinch
$4\sigma$. Every stable class of $M_4$ pinches to
$4j\sigma\in\ker(4:\mathbb Z/16\{\sigma\}\to
\mathbb Z/16\{\sigma\})$. After subtracting
$j\widehat a$, it comes from the bottom $S^7$, whose
James--Hopf image is zero by \textup{(D.2.43i)}. Hence
\[
 JH_*:\pi_{15}^{s}M_4\longrightarrow\pi_{15}D_2M_4
 \quad\text{is zero}.
 \tag{D.2.43n}
\]

It remains to handle the additional stable quotient
classes in \textup{(D.2.43k)}. As in D.2.38,
\[
 \pi_{15}^{s}M(\mathbb Z/2,8)=\pi_7(\mathbb S/2),
 \qquad 2\pi_7(\mathbb S/2)=0.
 \tag{D.2.43o}
\]
Indeed $2\id_{\mathbb S/2}$ factors through the top
pinch followed by zero or $i\eta$; the pinch of a seventh
homotopy class is a multiple of $\nu^2$, and
$\eta\nu^2=0$. The low stable stems used here and below
are listed in \cite[Section~2, (2.1), p.~149]{I_Oka1986}.
The Moore exact sequence shows that this group is
generated by a bottom $\sigma$-class and a coextension
of $\nu^2$.

For that coextension use the actual class
\[
 t:S^{14}\longrightarrow M(\mathbb Z/2,7),
 \qquad qt=\nu_8^2,
 \tag{D.2.43p}
\]
provided by
\cite[Lemma~2.8(ii), p.~523]{EH5MorisugiMukai2000}.
Its suspension
$\xi=Et:S^{15}\to M(\mathbb Z/2,8)\subset B$
has the required stable pinch. The second actual factor
in \textup{(D.2.43j)} gives
\[
 \partial\xi=i\eta_7qt=i\eta_7\nu_8^2=0.
 \tag{D.2.43q}
\]
Here $\eta_7\nu_8\in\pi_{11}S^7$ is already in the
stable range, where the fourth stem is zero.
Thus $\xi$ has an actual lift to $\pi_{15}K$, whose
stable class has zero James--Hopf image. The other
Moore quotient generator is covered by the actual map
$r\sigma_8:S^{15}\to S^8_r\to K$, again with zero
James--Hopf image.

Finally take any $x\in\pi_{15}^{s}K$. If its quotient
coordinates are $(a,\xi')$, exactness of the stable
cofiber sequence gives $4a+\eta q\xi'=0$.
The second term vanishes, since $q\xi'$ is a multiple
of $\nu^2$. Thus $a=4j\sigma$. Subtract the old
Moore-four class, a bottom-$r$ sphere class, and the
actual lift from \textup{(D.2.43q)}, according to these
quotient coordinates. The remainder comes from the
bottom $S^7$. Each summand has zero James--Hopf image,
the last by \textup{(D.2.43i)}. We have proved
\[
 JH_*:\pi_{15}^{s}K\longrightarrow\pi_{15}D_2K
 \quad\text{is zero}.
 \tag{D.2.43r}
\]
For the stable bottom remainder, no separate
unstabilization claim is needed. Exactness of
\textup{(D.2.43g)} makes $\partial_K$ injective, and
$2c$ has order four, proving \textup{(D.2.43e)}.

The additional free $S^9_R$ contributes stable classes
through its bottom inclusion, but $D_2S^9$ starts in
dimension eighteen. Their James--Hopf images in degree
fifteen are therefore zero. Next, attaching cells of
dimensions ten and eleven induces a surjection on
fifteenth stable homotopy: the corresponding quotient
terms are sums of $\pi_5^s\mathbb S$ and
$\pi_4^s\mathbb S$, both zero. The quadratic target group
remains unchanged by the dimension argument above, so
naturality prevents any enlargement of the James--Hopf
image at these stages. A twelve-cell is the first
possibility, since $\pi_{15}^sS^{12}=\pi_3^s\mathbb S$.
This proves the final finite-range assertion while
retaining the unresolved full-target test
\textup{(D.2.43d)}.
\end{proof}
